\documentclass[11pt,a4paper]{article}

% Essential packages
\usepackage[utf8]{inputenc}
\usepackage[T1]{fontenc}
\usepackage{amsmath,amssymb,amsthm}
\usepackage{graphicx}
\usepackage{float}
\usepackage{booktabs}
\usepackage{hyperref}
\usepackage{xcolor}
\usepackage{pythontex}
\usepackage[margin=1in]{geometry}

% Custom commands
\newcommand{\dd}{\mathrm{d}}
\newcommand{\ee}{\mathrm{e}}
\newcommand{\vect}[1]{\mathbf{#1}}

% Document metadata
\title{Molecular Dynamics Simulation:\\From Lennard-Jones Potentials to Thermodynamic Properties}
\author{Computational Chemistry}
\date{\today}

\begin{document}

\maketitle

\begin{abstract}
Molecular dynamics (MD) simulations provide atomistic insight into the behavior of matter by solving Newton's equations of motion for systems of interacting particles. This document develops a complete MD simulation framework, starting from the Lennard-Jones potential for pairwise interactions, implementing the velocity Verlet integration algorithm, and extracting thermodynamic properties including temperature, pressure, and radial distribution functions. We simulate a simple Lennard-Jones fluid, demonstrating phase behavior and energy conservation.
\end{abstract}

\tableofcontents
\newpage

%------------------------------------------------------------------------------
\section{Introduction: Simulating Matter at the Atomic Scale}
%------------------------------------------------------------------------------

Molecular dynamics simulations bridge the gap between quantum mechanical descriptions of molecular interactions and macroscopic thermodynamic properties. By numerically solving Newton's equations of motion for a system of $N$ particles, we can:

\begin{itemize}
    \item Observe molecular-level dynamics in real time
    \item Calculate thermodynamic properties from first principles
    \item Study phase transitions and transport phenomena
    \item Investigate protein folding, chemical reactions, and material properties
\end{itemize}

The fundamental algorithm is straightforward:
\begin{enumerate}
    \item Define interatomic potentials that describe particle interactions
    \item Initialize positions and velocities
    \item Compute forces on each particle: $\vect{F}_i = -\nabla_i U$
    \item Integrate equations of motion: $\vect{F}_i = m_i \ddot{\vect{r}}_i$
    \item Repeat, collecting statistics along the trajectory
\end{enumerate}

\begin{pycode}
import numpy as np
import matplotlib.pyplot as plt
from mpl_toolkits.mplot3d import Axes3D

# Set random seed for reproducibility
np.random.seed(42)

# Style configuration
plt.rcParams['figure.figsize'] = (10, 6)
plt.rcParams['font.size'] = 11
plt.rcParams['axes.labelsize'] = 12
plt.rcParams['axes.titlesize'] = 14
\end{pycode}

%------------------------------------------------------------------------------
\section{The Lennard-Jones Potential}
%------------------------------------------------------------------------------

The Lennard-Jones (LJ) potential is the most widely used model for van der Waals interactions between neutral atoms:

\begin{equation}
U_{\text{LJ}}(r) = 4\varepsilon \left[\left(\frac{\sigma}{r}\right)^{12} - \left(\frac{\sigma}{r}\right)^{6}\right]
\label{eq:lennard-jones}
\end{equation}

where:
\begin{itemize}
    \item $\varepsilon$ is the depth of the potential well (energy scale)
    \item $\sigma$ is the distance at which $U = 0$ (length scale)
    \item $r$ is the interparticle distance
\end{itemize}

The $r^{-12}$ term models Pauli repulsion at short range, while the $r^{-6}$ term models the attractive London dispersion force. The potential has a minimum at $r_{\text{min}} = 2^{1/6}\sigma \approx 1.122\sigma$ with depth $-\varepsilon$.

\subsection{Force Derivation}

The force between particles is the negative gradient of the potential:

\begin{equation}
F_{\text{LJ}}(r) = -\frac{\dd U_{\text{LJ}}}{\dd r} = \frac{24\varepsilon}{r}\left[2\left(\frac{\sigma}{r}\right)^{12} - \left(\frac{\sigma}{r}\right)^{6}\right]
\label{eq:lj-force}
\end{equation}

In vector form, the force on particle $i$ due to particle $j$ is:

\begin{equation}
\vect{F}_{ij} = F_{\text{LJ}}(r_{ij}) \frac{\vect{r}_{ij}}{r_{ij}}
\end{equation}

\begin{pycode}
# Lennard-Jones potential and force
def lennard_jones_potential(r, epsilon=1.0, sigma=1.0):
    """Lennard-Jones pair potential."""
    sr6 = (sigma / r)**6
    sr12 = sr6**2
    return 4 * epsilon * (sr12 - sr6)

def lennard_jones_force(r, epsilon=1.0, sigma=1.0):
    """Magnitude of Lennard-Jones force."""
    sr6 = (sigma / r)**6
    sr12 = sr6**2
    return 24 * epsilon / r * (2 * sr12 - sr6)

# Plot potential and force
fig, axes = plt.subplots(1, 2, figsize=(12, 5))

r = np.linspace(0.9, 3.0, 500)
U = lennard_jones_potential(r)
F = lennard_jones_force(r)

# Left: Potential
ax = axes[0]
ax.plot(r, U, 'b-', linewidth=2)
ax.axhline(0, color='gray', linestyle='--', alpha=0.5)
ax.axvline(2**(1/6), color='red', linestyle=':', alpha=0.7, label=r'$r_{min} = 2^{1/6}\sigma$')

# Mark key points
r_min = 2**(1/6)
ax.plot(r_min, -1, 'ro', markersize=8)
ax.annotate(r'$(-\varepsilon)$', xy=(r_min, -1), xytext=(r_min + 0.2, -0.7),
            fontsize=11, color='red')
ax.plot(1.0, 0, 'go', markersize=8)
ax.annotate(r'$\sigma$', xy=(1.0, 0), xytext=(0.85, 0.5), fontsize=11, color='green')

ax.set_xlabel(r'Distance $r/\sigma$')
ax.set_ylabel(r'Potential $U/\varepsilon$')
ax.set_title('Lennard-Jones Potential')
ax.set_xlim(0.9, 3.0)
ax.set_ylim(-1.5, 2.0)
ax.legend()
ax.grid(True, alpha=0.3)

# Right: Force
ax = axes[1]
ax.plot(r, F, 'r-', linewidth=2)
ax.axhline(0, color='gray', linestyle='--', alpha=0.5)
ax.axvline(2**(1/6), color='blue', linestyle=':', alpha=0.7)

ax.fill_between(r[F > 0], 0, F[F > 0], alpha=0.3, color='red', label='Repulsive')
ax.fill_between(r[F < 0], 0, F[F < 0], alpha=0.3, color='blue', label='Attractive')

ax.set_xlabel(r'Distance $r/\sigma$')
ax.set_ylabel(r'Force $F\sigma/\varepsilon$')
ax.set_title('Lennard-Jones Force')
ax.set_xlim(0.9, 3.0)
ax.set_ylim(-3, 5)
ax.legend()
ax.grid(True, alpha=0.3)

plt.tight_layout()
plt.savefig('molecular_dynamics_potential.pdf', bbox_inches='tight')
plt.close()
\end{pycode}

\begin{figure}[H]
\centering
\includegraphics[width=\textwidth]{molecular_dynamics_potential.pdf}
\caption{The Lennard-Jones potential and its derivative. Left: The potential energy curve shows strong repulsion at short range ($r < \sigma$), a minimum at $r_{\text{min}} = 2^{1/6}\sigma$ with depth $-\varepsilon$, and asymptotic approach to zero at large distances. Right: The force is repulsive (positive) for $r < r_{\text{min}}$ and attractive (negative) for $r > r_{\text{min}}$. At the equilibrium distance $r_{\text{min}}$, the force is zero.}
\label{fig:lj-potential}
\end{figure}

%------------------------------------------------------------------------------
\section{Reduced Units}
%------------------------------------------------------------------------------

MD simulations typically use \textit{reduced units} where $\varepsilon$, $\sigma$, and the particle mass $m$ are set to unity. All quantities are then expressed in terms of these fundamental scales:

\begin{table}[H]
\centering
\begin{tabular}{lcc}
\toprule
Quantity & Reduced Unit & For Argon \\
\midrule
Length & $\sigma$ & 3.4 \AA \\
Energy & $\varepsilon$ & $1.65 \times 10^{-21}$ J \\
Mass & $m$ & $6.63 \times 10^{-26}$ kg \\
Time & $\sigma\sqrt{m/\varepsilon}$ & 2.15 ps \\
Temperature & $\varepsilon/k_B$ & 120 K \\
Pressure & $\varepsilon/\sigma^3$ & 42 MPa \\
\bottomrule
\end{tabular}
\caption{Reduced units and their values for liquid argon.}
\label{tab:units}
\end{table}

%------------------------------------------------------------------------------
\section{The Velocity Verlet Algorithm}
%------------------------------------------------------------------------------

The velocity Verlet algorithm is the most popular integrator for MD due to its time-reversibility, symplectic nature, and good energy conservation:

\begin{align}
\vect{r}(t + \Delta t) &= \vect{r}(t) + \vect{v}(t)\Delta t + \frac{1}{2}\vect{a}(t)\Delta t^2 \label{eq:verlet-r}\\
\vect{v}(t + \Delta t) &= \vect{v}(t) + \frac{1}{2}[\vect{a}(t) + \vect{a}(t + \Delta t)]\Delta t \label{eq:verlet-v}
\end{align}

where $\vect{a} = \vect{F}/m$. The algorithm proceeds as:
\begin{enumerate}
    \item Calculate forces at time $t$
    \item Update positions using Eq.~\eqref{eq:verlet-r}
    \item Calculate forces at time $t + \Delta t$
    \item Update velocities using Eq.~\eqref{eq:verlet-v}
\end{enumerate}

\begin{pycode}
class LennardJonesSimulation:
    """
    Molecular dynamics simulation of a Lennard-Jones fluid.

    Uses reduced units: epsilon = sigma = mass = 1
    """

    def __init__(self, n_particles, box_length, temperature, dt=0.002, cutoff=2.5):
        """
        Initialize the simulation.

        Parameters:
        -----------
        n_particles : int
            Number of particles
        box_length : float
            Side length of cubic simulation box
        temperature : float
            Initial temperature (reduced units)
        dt : float
            Time step (reduced units)
        cutoff : float
            Potential cutoff distance
        """
        self.n = n_particles
        self.L = box_length
        self.T = temperature
        self.dt = dt
        self.rc = cutoff
        self.rc2 = cutoff**2

        # Potential shift for continuity at cutoff
        self.U_shift = 4 * ((1/cutoff)**12 - (1/cutoff)**6)

        # Initialize positions on a lattice
        self.positions = self._init_lattice()

        # Initialize velocities from Maxwell-Boltzmann distribution
        self.velocities = self._init_velocities()

        # Force array
        self.forces = np.zeros_like(self.positions)

        # Observables
        self.potential_energy = 0
        self.kinetic_energy = 0
        self.virial = 0

    def _init_lattice(self):
        """Place particles on a simple cubic lattice."""
        n_side = int(np.ceil(self.n**(1/3)))
        spacing = self.L / n_side
        positions = []

        for i in range(n_side):
            for j in range(n_side):
                for k in range(n_side):
                    if len(positions) < self.n:
                        positions.append([
                            (i + 0.5) * spacing,
                            (j + 0.5) * spacing,
                            (k + 0.5) * spacing
                        ])

        return np.array(positions)

    def _init_velocities(self):
        """Initialize velocities from Maxwell-Boltzmann distribution."""
        velocities = np.random.randn(self.n, 3) * np.sqrt(self.T)

        # Remove center of mass motion
        velocities -= np.mean(velocities, axis=0)

        # Scale to exact temperature
        ke = 0.5 * np.sum(velocities**2)
        current_T = 2 * ke / (3 * self.n)
        velocities *= np.sqrt(self.T / current_T)

        return velocities

    def _minimum_image(self, dr):
        """Apply minimum image convention for periodic boundaries."""
        return dr - self.L * np.round(dr / self.L)

    def compute_forces(self):
        """
        Compute forces, potential energy, and virial.

        Uses the minimum image convention for periodic boundaries.
        """
        self.forces.fill(0)
        self.potential_energy = 0
        self.virial = 0

        for i in range(self.n - 1):
            for j in range(i + 1, self.n):
                dr = self.positions[j] - self.positions[i]
                dr = self._minimum_image(dr)
                r2 = np.sum(dr**2)

                if r2 < self.rc2:
                    r2_inv = 1 / r2
                    r6_inv = r2_inv**3
                    r12_inv = r6_inv**2

                    # Potential (shifted)
                    U = 4 * (r12_inv - r6_inv) - self.U_shift
                    self.potential_energy += U

                    # Force magnitude / r
                    f_over_r = 24 * r2_inv * (2 * r12_inv - r6_inv)

                    # Force vector
                    f = f_over_r * dr
                    self.forces[i] -= f
                    self.forces[j] += f

                    # Virial for pressure calculation
                    self.virial += f_over_r * r2

    def velocity_verlet_step(self):
        """Perform one velocity Verlet integration step."""
        # Half-step velocity update
        self.velocities += 0.5 * self.forces * self.dt

        # Full position update
        self.positions += self.velocities * self.dt

        # Apply periodic boundary conditions
        self.positions = self.positions % self.L

        # Compute new forces
        self.compute_forces()

        # Complete velocity update
        self.velocities += 0.5 * self.forces * self.dt

        # Calculate kinetic energy
        self.kinetic_energy = 0.5 * np.sum(self.velocities**2)

    def get_temperature(self):
        """Calculate instantaneous temperature."""
        return 2 * self.kinetic_energy / (3 * self.n)

    def get_pressure(self):
        """Calculate pressure using virial theorem."""
        density = self.n / self.L**3
        T = self.get_temperature()
        return density * T + self.virial / (3 * self.L**3)

    def get_total_energy(self):
        """Calculate total energy."""
        return self.kinetic_energy + self.potential_energy

    def apply_thermostat(self, target_T):
        """Velocity rescaling thermostat."""
        current_T = self.get_temperature()
        if current_T > 0:
            scale = np.sqrt(target_T / current_T)
            self.velocities *= scale

# Run a short simulation
np.random.seed(42)
n_particles = 108  # 4x4x4 FCC-like arrangement
density = 0.8  # Reduced density (liquid state)
box_length = (n_particles / density)**(1/3)
temperature = 1.0  # Reduced temperature

sim = LennardJonesSimulation(n_particles, box_length, temperature)
sim.compute_forces()

# Equilibration with thermostat
n_equil = 500
for _ in range(n_equil):
    sim.velocity_verlet_step()
    if _ % 50 == 0:
        sim.apply_thermostat(temperature)

# Production run (NVE - constant energy)
n_steps = 2000
energies = {'kinetic': [], 'potential': [], 'total': []}
temperatures = []
pressures = []

for step in range(n_steps):
    sim.velocity_verlet_step()

    energies['kinetic'].append(sim.kinetic_energy / n_particles)
    energies['potential'].append(sim.potential_energy / n_particles)
    energies['total'].append(sim.get_total_energy() / n_particles)
    temperatures.append(sim.get_temperature())
    pressures.append(sim.get_pressure())

# Convert to arrays
for key in energies:
    energies[key] = np.array(energies[key])
temperatures = np.array(temperatures)
pressures = np.array(pressures)
time_array = np.arange(n_steps) * sim.dt
\end{pycode}

%------------------------------------------------------------------------------
\section{Simulation Results}
%------------------------------------------------------------------------------

\subsection{Energy Conservation}

A hallmark of a good MD integrator is energy conservation in the microcanonical (NVE) ensemble. The velocity Verlet algorithm achieves excellent long-term energy conservation.

\begin{pycode}
# Energy conservation plot
fig, axes = plt.subplots(2, 2, figsize=(12, 10))

# Top left: Energy components
ax = axes[0, 0]
ax.plot(time_array, energies['kinetic'], 'r-', alpha=0.7, label='Kinetic')
ax.plot(time_array, energies['potential'], 'b-', alpha=0.7, label='Potential')
ax.plot(time_array, energies['total'], 'k-', linewidth=2, label='Total')

ax.set_xlabel('Time (reduced units)')
ax.set_ylabel('Energy per particle')
ax.set_title('Energy Components')
ax.legend()
ax.grid(True, alpha=0.3)

# Top right: Energy conservation detail
ax = axes[0, 1]
E_total = energies['total']
E_mean = np.mean(E_total)
E_drift = (E_total - E_total[0]) * 100 / np.abs(E_mean)  # Percentage drift

ax.plot(time_array, E_drift, 'k-', linewidth=1)
ax.axhline(0, color='gray', linestyle='--', alpha=0.5)

ax.set_xlabel('Time (reduced units)')
ax.set_ylabel('Energy drift (%)')
ax.set_title(f'Energy Conservation (drift: {np.std(E_drift):.4f}%)')
ax.grid(True, alpha=0.3)

# Bottom left: Temperature fluctuations
ax = axes[1, 0]
ax.plot(time_array, temperatures, 'r-', alpha=0.7)
ax.axhline(np.mean(temperatures), color='k', linestyle='--',
           label=f'Mean: {np.mean(temperatures):.3f}')
ax.fill_between(time_array,
                np.mean(temperatures) - np.std(temperatures),
                np.mean(temperatures) + np.std(temperatures),
                alpha=0.2, color='red')

ax.set_xlabel('Time (reduced units)')
ax.set_ylabel('Temperature')
ax.set_title(f'Temperature (std: {np.std(temperatures):.3f})')
ax.legend()
ax.grid(True, alpha=0.3)

# Bottom right: Pressure fluctuations
ax = axes[1, 1]
ax.plot(time_array, pressures, 'b-', alpha=0.7)
ax.axhline(np.mean(pressures), color='k', linestyle='--',
           label=f'Mean: {np.mean(pressures):.3f}')

ax.set_xlabel('Time (reduced units)')
ax.set_ylabel('Pressure')
ax.set_title(f'Pressure (std: {np.std(pressures):.3f})')
ax.legend()
ax.grid(True, alpha=0.3)

plt.tight_layout()
plt.savefig('molecular_dynamics_energy.pdf', bbox_inches='tight')
plt.close()
\end{pycode}

\begin{figure}[H]
\centering
\includegraphics[width=\textwidth]{molecular_dynamics_energy.pdf}
\caption{Thermodynamic properties during MD simulation. Top left: Energy components showing the exchange between kinetic and potential energy while total energy remains constant. Top right: Energy drift as a percentage of mean energy, demonstrating excellent conservation by the velocity Verlet integrator. Bottom left: Temperature fluctuations around the mean---these are expected in the NVE ensemble and decrease as $1/\sqrt{N}$. Bottom right: Pressure fluctuations; the mean pressure includes both ideal gas and virial contributions.}
\label{fig:energy}
\end{figure}

\subsection{Radial Distribution Function}

The radial distribution function (RDF) $g(r)$ measures the probability of finding a particle at distance $r$ from another particle, relative to an ideal gas. It provides structural information about the fluid:

\begin{equation}
g(r) = \frac{V}{N^2} \left\langle \sum_i \sum_{j \neq i} \delta(r - r_{ij}) \right\rangle
\end{equation}

\begin{pycode}
def compute_rdf(positions, box_length, n_bins=100, r_max=None):
    """
    Compute the radial distribution function.

    Parameters:
    -----------
    positions : ndarray
        Particle positions (N x 3)
    box_length : float
        Simulation box size
    n_bins : int
        Number of histogram bins
    r_max : float
        Maximum distance (default: half box length)
    """
    if r_max is None:
        r_max = box_length / 2

    n = len(positions)
    dr = r_max / n_bins
    hist = np.zeros(n_bins)

    for i in range(n - 1):
        for j in range(i + 1, n):
            dpos = positions[j] - positions[i]
            dpos = dpos - box_length * np.round(dpos / box_length)
            r = np.sqrt(np.sum(dpos**2))

            if r < r_max:
                bin_idx = int(r / dr)
                hist[bin_idx] += 2  # Count both i-j and j-i

    # Normalize
    r_edges = np.linspace(0, r_max, n_bins + 1)
    r_centers = 0.5 * (r_edges[:-1] + r_edges[1:])
    shell_volumes = 4 * np.pi * r_centers**2 * dr

    density = n / box_length**3
    g_r = hist / (n * density * shell_volumes)

    return r_centers, g_r

# Collect RDF samples
rdf_samples = []
for _ in range(200):
    sim.velocity_verlet_step()
    if _ % 10 == 0:
        r_vals, g_vals = compute_rdf(sim.positions, sim.L)
        rdf_samples.append(g_vals)

g_mean = np.mean(rdf_samples, axis=0)
g_std = np.std(rdf_samples, axis=0)

# Plot RDF
fig, ax = plt.subplots(figsize=(10, 6))

ax.plot(r_vals, g_mean, 'b-', linewidth=2, label='Simulation')
ax.fill_between(r_vals, g_mean - g_std, g_mean + g_std, alpha=0.3, color='blue')
ax.axhline(1, color='gray', linestyle='--', alpha=0.7, label='Ideal gas')

# Mark key features
peaks = [1.0, 1.95]  # Approximate peak positions for LJ liquid
for i, p in enumerate(peaks):
    idx = np.argmin(np.abs(r_vals - p))
    if g_mean[idx] > 1:
        ax.annotate(f'{i+1}st shell', xy=(r_vals[idx], g_mean[idx]),
                    xytext=(r_vals[idx] + 0.2, g_mean[idx] + 0.5),
                    arrowprops=dict(arrowstyle='->', color='red'),
                    fontsize=10, color='red')

ax.set_xlabel(r'Distance $r/\sigma$')
ax.set_ylabel(r'$g(r)$')
ax.set_title(f'Radial Distribution Function ($\\rho^* = {density:.1f}$, $T^* = {temperature:.1f}$)')
ax.legend()
ax.grid(True, alpha=0.3)
ax.set_xlim(0, sim.L / 2)
ax.set_ylim(0, 3.5)

plt.tight_layout()
plt.savefig('molecular_dynamics_rdf.pdf', bbox_inches='tight')
plt.close()
\end{pycode}

\begin{figure}[H]
\centering
\includegraphics[width=0.9\textwidth]{molecular_dynamics_rdf.pdf}
\caption{Radial distribution function $g(r)$ for the Lennard-Jones fluid at liquid density ($\rho^* = 0.8$) and temperature ($T^* = 1.0$). The first peak near $r = \sigma$ corresponds to the first coordination shell of nearest neighbors. The second peak indicates the second shell. At large distances, $g(r) \to 1$ (ideal gas behavior). The oscillatory structure is characteristic of liquids, contrasting with the delta-function peaks of solids and the uniform $g(r) = 1$ of gases.}
\label{fig:rdf}
\end{figure}

\subsection{Phase Diagram Exploration}

The Lennard-Jones system exhibits gas, liquid, and solid phases. The phase depends on reduced density $\rho^* = \rho\sigma^3$ and reduced temperature $T^* = k_BT/\varepsilon$.

\begin{pycode}
# Scan different state points
state_points = [
    (0.1, 2.0, 'Gas'),
    (0.8, 1.0, 'Liquid'),
    (1.0, 0.5, 'Solid'),
]

fig, axes = plt.subplots(1, 3, figsize=(15, 5))

for ax, (density, temp, phase) in zip(axes, state_points):
    # Quick simulation at this state point
    np.random.seed(42)
    n_test = 64
    L_test = (n_test / density)**(1/3)

    sim_test = LennardJonesSimulation(n_test, L_test, temp)
    sim_test.compute_forces()

    # Short equilibration
    for _ in range(300):
        sim_test.velocity_verlet_step()
        if _ % 20 == 0:
            sim_test.apply_thermostat(temp)

    # Plot 3D snapshot
    pos = sim_test.positions
    ax.scatter(pos[:, 0], pos[:, 1], c=pos[:, 2], cmap='viridis',
               s=500, alpha=0.7, edgecolors='white')

    ax.set_xlim(0, L_test)
    ax.set_ylim(0, L_test)
    ax.set_xlabel('x')
    ax.set_ylabel('y')
    ax.set_title(f'{phase}\n($\\rho^* = {density}$, $T^* = {temp}$)')
    ax.set_aspect('equal')
    ax.grid(True, alpha=0.3)

plt.tight_layout()
plt.savefig('molecular_dynamics_phases.pdf', bbox_inches='tight')
plt.close()
\end{pycode}

\begin{figure}[H]
\centering
\includegraphics[width=\textwidth]{molecular_dynamics_phases.pdf}
\caption{Snapshots of the Lennard-Jones system in different phases (2D projection, colored by z-coordinate). Left: Gas phase at low density---particles are well-separated with no persistent structure. Center: Liquid phase at intermediate density---particles form a disordered but cohesive structure with short-range order. Right: Solid phase at high density and low temperature---particles arrange in an ordered crystalline lattice. These phases emerge naturally from the interplay between kinetic energy (temperature) and potential energy (density).}
\label{fig:phases}
\end{figure}

%------------------------------------------------------------------------------
\section{Conclusions}
%------------------------------------------------------------------------------

This document has presented a complete molecular dynamics simulation framework:

\begin{enumerate}
    \item \textbf{Lennard-Jones potential}: The $12$-$6$ form captures essential physics of neutral atom interactions, with repulsion at short range and attraction at longer distances.

    \item \textbf{Velocity Verlet integration}: This symplectic integrator provides excellent energy conservation, enabling accurate NVE simulations over long timescales.

    \item \textbf{Thermodynamic properties}: Temperature, pressure, and energy are directly calculable from particle positions and velocities. The virial theorem connects microscopic forces to macroscopic pressure.

    \item \textbf{Structural analysis}: The radial distribution function reveals the characteristic short-range order of liquids, distinguishing them from gases and solids.

    \item \textbf{Phase behavior}: By varying density and temperature, the simulation captures gas, liquid, and solid phases---all emerging from the same simple interatomic potential.
\end{enumerate}

The simulated system of \pyc{print(n_particles)} particles showed:
\begin{itemize}
    \item Energy conservation to \pyc{print(f'{np.std(E_drift):.4f}')}{\%}
    \item Mean temperature: \pyc{print(f'{np.mean(temperatures):.3f}')} (target: \pyc{print(f'{temperature:.1f}')})
    \item Mean pressure: \pyc{print(f'{np.mean(pressures):.3f}')} (reduced units)
    \item First RDF peak at $r \approx \sigma$, characteristic of liquid structure
\end{itemize}

Extensions of this framework include more realistic potentials (EAM for metals, force fields for biomolecules), constant temperature and pressure ensembles, long-range electrostatics (Ewald summation), and reactive MD for chemical transformations.

%------------------------------------------------------------------------------
\section{References}
%------------------------------------------------------------------------------

\begin{enumerate}
    \item Allen, M.P. \& Tildesley, D.J. (2017). \textit{Computer Simulation of Liquids}, 2nd ed. Oxford University Press.

    \item Frenkel, D. \& Smit, B. (2002). \textit{Understanding Molecular Simulation}, 2nd ed. Academic Press.

    \item Rapaport, D.C. (2004). \textit{The Art of Molecular Dynamics Simulation}, 2nd ed. Cambridge University Press.

    \item Lennard-Jones, J.E. (1924). On the Determination of Molecular Fields. \textit{Proc. R. Soc. Lond. A}, 106(738), 463-477.

    \item Verlet, L. (1967). Computer ``Experiments'' on Classical Fluids. \textit{Physical Review}, 159(1), 98-103.

    \item Hansen, J.P. \& Verlet, L. (1969). Phase Transitions of the Lennard-Jones System. \textit{Physical Review}, 184(1), 151-161.
\end{enumerate}

\end{document}
